\documentclass[9pt,twocolumn]{extarticle}
\usepackage{kiee-review}
% Same typesetting in marked and clean copies; only the revision color differs.
\ifdefined\CleanCopy
  \definecolor{revisionblue}{RGB}{0,0,0}
\else
  \definecolor{revisionblue}{RGB}{0,70,180}
\fi
\newenvironment{revision}{\par\begingroup\color{revisionblue}}{\par\endgroup}
\newcommand{\rev}[1]{{\color{revisionblue}#1}}
\setlength{\emergencystretch}{2em}
\captionsetup{font={footnotesize,color=revisionblue},labelfont={bf,color=revisionblue}}
% Mark added references while preserving unchanged submitted bibliography text.
\makeatletter
\g@addto@macro\UrlBreaks{\do\-}
\@for\bibkey:=leucker2009runtime,bauer2011runtime,uppaalofficial,uppaalname,uppaalqueries\do{%
  \expandafter\gdef\csname changedbib@\bibkey\endcsname{1}}
\let\originalbibitem\@bibitem
\def\@bibitem#1{\ifcsname changedbib@#1\endcsname\color{revisionblue}\else\color{black}\fi\originalbibitem{#1}}
\makeatother


\kieeenglishtitle{\rev{A Feasibility Framework for Temporal and Conditional Consistency Checking of Sequential Chain-of-Causation Annotations}}
\kieekoreantitle{\rev{연속 인과관계(CoC) 주석의 시간·조건 일관성 검사 프레임워크: 적용 가능성과 근거 추적}}

\begin{document}
\makekieefrontmatter{%
\begin{kieeabstract}
\rev{Sequential Chain-of-Causation (CoC) annotations may omit release evidence or justify an action that conflicts with an earlier obligation. We link obligations to condition-specific evidence and use UPPAAL to distinguish confirmed conflicts from insufficient evidence and record the first issue event, obligation, and reason. This study evaluates feasibility and traceability under authored symbolic inputs. The natural-data inventory contains 309 adjacent pairs from 90 scenes; four reviewers assessed a selected subset of 27 scene groups, with textual-consistency Fleiss' kappa 0.027301 and no recorded consensus rationale. In 144 authored scenarios (72 contradictions, 36 consistent cases, and 36 insufficient-evidence cases), the frozen legacy structured path matched 96 decisions and a simple finite-state comparator matched all 144. The automatic text path returned UNKNOWN for every case. After analyzing 48 legacy mismatches, a separate condition-indexed UPPAAL model matched 144/144 decisions, 72/72 first violation events, and 108/108 issue locations, obligations, and reasons on the same cases. This is a post-development same-case retest, not independent validation. Python--UPPAAL agreement checks implementation consistency; neither natural-error accuracy, automated semantic extraction, accuracy superiority of UPPAAL, nor vehicle safety is established.}
\end{kieeabstract}
\kieekeywords{\rev{Chain-of-Causation, Temporal consistency, Condition-linked obligations, Model checking, Feasibility, Evidence traceability}}
}

\section{서론}
\label{sec:intro}

\begin{revision}

비전--언어 주행 모델은 장면을 해석하고 행동 이유를 언어로 설명한다 \cite{sima2024drivelm,tian2024drivevlm}. Chain-of-Causation(CoC)은 이런 행동 이유를 인과관계로 연결한 주석이다. 하지만 이 방식은 같은 장면의 CoC들이 개별적으로는 자연스러워도 시간순으로 함께 읽으면 양립하지 않을 수 있다는 문제점이 있다. 생성ㆍ작성 과정에서 조건이 누락되거나, 앞서 제시한 요구와 양립하지 않는 행동이 기록될 수 있기 때문이다. 이러한 불일치 때문에 학습 자료의 행동 요구가 모호해지고, 같은 자료를 다시 검사했을 때 일관된 결과를 얻기 어렵다.

다음은 연구 동기를 설명하는 가상 예시이며 자연 자료에서 발견한 오류가 아니다. 앞선 CoC가 ``보행자가 통과할 때까지 정지한다''고 하고, 후속 CoC가 ``신호가 녹색으로 바뀌었으므로 출발한다''고 했다고 하자. 녹색 신호는 출발의 새로운 근거이지만 보행자가 통과했다는 증거는 아니다. 보행자가 여전히 횡단 중임이 확인되면 출발은 이전 정지 요구와 충돌한다. 통과 여부가 제시되지 않았다면 확정 모순이 아니라 정지 요구의 해제 근거가 부족한 경우이다. 보행자의 통과가 확인된 뒤 출발하면 해당 요구와 양립한다. 새로운 행동 근거 B가 있다는 이유만으로 이전 요구 A를 지워서는 안 된다.

본 연구는 이러한 시간 순서와 조건의 관계를 UPPAAL 모델로 검사하고 문제 사건과 근거를 추적한다. UPPAAL은 Uppsala 대학과 Aalborg 대학이 공동 개발한 시간 오토마타 기반 모델 검사 도구의 이름이며, 두 대학 이름을 연결한 명칭이다 \cite{uppaalofficial,uppaalname,behrmann2004uppaal}. 문장의 요구를 구조화하는 과정과 그 조건의 참ㆍ거짓을 확인하는 과정은 별개다. 현재 자동 변환기는 제한된 영어 행동ㆍ조건 문구를 찾지만 관측 증거를 확정하지 않는다. 따라서 자동 텍스트 경로와 설계자가 작성한 조건별 프레임 경로를 구분한다.

\begin{figure*}[t]
\centering
\color{revisionblue}
\begin{tikzpicture}[
 box/.style={draw=revisionblue,rounded corners,align=center,text width=3.4cm,minimum height=11mm,font=\footnotesize,text=revisionblue},
 flow/.style={-{Latex},draw=revisionblue,thick}]
\node[box] (text) at (0,0) {Stored CoC text\\rule-based phrase extraction\\observed evidence unresolved};
\node[box] (frame) at (0,-1.7) {Authored text + frames\\condition IDs and source spans\\fixed scenario expectations};
\node[box] (check) at (5,-.85) {Frozen single-obligation path\\FSM comparator\\condition-indexed UPPAAL};
\node[box] (trace) at (10,-.85) {Decision + first issue event\\obligation + reason\\XML / queries / logs};
\draw[flow] (text.east) -- (check.north west);
\draw[flow] (frame.east) -- (check.south west);
\draw[flow] (check.east) -- (trace.west);
\end{tikzpicture}

\bifigcaption{자동 문구 추출과 작성 프레임을 구분한 검사 흐름. 원문 구절, 조건별 의무, 최초 문제 사건과 사유를 연결한다.}{Separate automatic-text and authored-frame paths with links from source spans to obligations and issue records.}
\label{fig:overview-pipeline}
\end{figure*}

연구 질문은 ``연속 CoC의 의무와 해제 근거를 연결하여 행동 전환의 모순과 근거 부족을 구분하고, 최초 문제 사건ㆍ관련 의무ㆍ사유를 추적할 수 있는가?''이다. 이를 확인하기 위해 기존 단일 의무 검사기의 최초 결과를 보존하고, 이전 사건 1/3/5개를 보는 검사와 별도 유한상태기계(FSM)를 비교했다. 최초 실패 분석 후 조건별 UPPAAL 모델을 개발하여 같은 사례를 재시험했다. 이 재시험을 독립 성능 검증으로 해석하지 않는다.

기여는 세 가지다. 첫째, 원문 구절--의무 ID--해제 조건--후속 증거--행동을 연결하고 모순과 근거 부족을 구분하는 표현을 제시한다. 둘째, 조건별 상태와 속성을 UPPAAL로 실행하고 최초 문제 위치와 사유를 추적하는 절차를 구체화한다. 셋째, 최초 통제 시험, 실패 분석, 동일 사례 재시험과 자연 자료의 해석 한계를 구분하여 적용 가능성을 보고한다. 단순 FSM도 같은 사례를 모두 판정했으므로 UPPAAL의 정확도 우월성을 주장하지 않는다.

중심 검사는 고정된 사건 순서와 조건을 다룬다. 보조 반응 시간 분석은 CoC의 행동 요구와 기록된 자차 속도 변화의 관계를 따로 살핀다. 정량적 제한 시간, 차량 동역학과 폐루프 제어는 중심 조건별 모델에 없다. 검증 범위는 명시된 의무와 증거에 한정하며, 언급되지 않은 위험의 완전성, 자연 오류 정확도, 실제 안전 또는 학습 성능 향상을 입증하지 않는다.
\end{revision}

\section{관련 연구}
\label{sec:related}

\begin{revision}

DriveLM은 지각ㆍ예측ㆍ계획의 질의응답을 그래프로 연결하고, DriveVLM은 장면 설명ㆍ분석ㆍ계층적 계획을 결합한다 \cite{sima2024drivelm,tian2024drivevlm}. 이들은 설명과 행동 생성의 관계를 다룬다. 본 연구는 저장된 설명을 다시 읽어 선행 의무와 후속 근거를 연결한다. 생성 모델의 계획 성능이나 설명의 실제 인과적 충실성을 재평가하는 방법은 아니다.

설명 검증에 가까운 최근 프리프린트는 언어 근거와 장면ㆍ궤적의 충실성 및 시각 교란에 대한 반응을 평가한다 \cite{mayumu2026faithfulness}. VLADriveBench는 언급ㆍ환각ㆍ모순ㆍ행동 정렬의 관측 지표와 CoT 개입을 함께 사용한다 \cite{nguyen2026vladrivebench}. 이런 검사는 생성된 설명과 행동이 맞는지 또는 설명을 바꾸었을 때 행동이 달라지는지를 묻는다. 본 논문은 모델을 다시 실행하거나 개입하지 않고 저장된 연속 설명의 의무와 해제 근거를 연결한다. 따라서 현재 출력은 실제 원인에 충실하다는 인과 검증의 대용이 아니다.

실행 중 검증(runtime verification)은 관측된 실행 경로가 주어진 명세를 만족하는지 검사한다 \cite{leucker2009runtime}. LTL/TLTL 모니터링은 유한 접두 경로에서 참ㆍ거짓ㆍ미결정을 구분하고 최초로 확정되는 시점을 다룬다 \cite{bauer2011runtime}. 따라서 이전 상태의 기억, 세 값의 판정과 최초 위반 추적 자체는 기존 기술이다. 본 논문의 미상은 명시된 해제 증거의 부족을 뜻하며, 미래 연장 전체에 대한 LTL 세 값 의미론을 구현했다는 뜻은 아니다. 명세로부터 범용 모니터를 자동 합성하는 것도 현재 범위에 없다.

UPPAAL은 시간 오토마타의 상태와 도달 속성을 검사한다 \cite{behrmann2004uppaal,uppaalqueries}. 본 연구의 중심 모델은 사건 배열을 순서대로 처리하는 이산 관찰기이며, 시간 오토마타 도구를 사용했다는 이유만으로 물리적 시간 제약을 검증한 것은 아니다. 조건 ID에 연결된 활성 의무와 증거를 상태 변수로 두고, 모순 및 증거 부족 상태의 도달 여부를 확인한다.

Scenic과 VerifAI는 자율 시스템 시험을 위한 상황 기술ㆍ생성 및 분석 도구를 제공한다 \cite{fremont2019scenic,dreossi2019verifai}. 이들의 환경 탐색ㆍ시뮬레이션과 달리 본 실험은 작성된 사건 배열을 고정해 검사한다. 본 연구의 통제 시나리오는 이 도구들과의 성능 경쟁 자료가 아니며 환경 다양성이나 실제 운전 안전을 대신 평가하지 않는다.

\begin{table*}[t]
\centering
\color{revisionblue}
\bitablecaption{관련 방법의 입력ㆍ표현ㆍ출력과 본 논문의 범위.}{Inputs, representations, outputs, and scope relative to prior methods.}
\label{tab:related-position}
\scriptsize
\begin{tabularx}{\textwidth}{L{.21\textwidth}L{.24\textwidth}L{.24\textwidth}Y}
\toprule
방법 & 입력과 표현 & 출력 & 본 연구와의 관계 \\
\midrule
DriveLM/DriveVLM \cite{sima2024drivelm,tian2024drivevlm} & 장면 관측, 언어 추론과 계획 & 설명ㆍ행동/계획 & 저장된 설명 사이의 의무 해제 연결은 본 연구의 검사 대상 \\
설명 충실성/CoT--행동 검증 \cite{mayumu2026faithfulness,nguyen2026vladrivebench} & 생성 설명ㆍ궤적, 관측 및 입력/CoT 개입 & 장면ㆍ행동 정렬, 개입 반응 & 본 논문은 저장된 주석의 조건 연결 검사; 생성 모델의 인과 충실성은 미검증 \\
실행 중 검증 \cite{leucker2009runtime} & 명세와 관측 실행 경로 & 경로의 명세 만족 여부 & 경로 검사 개념을 재사용; 자연어 해석 정답을 제공하지 않음 \\
LTL/TLTL 모니터 \cite{bauer2011runtime} & 논리식, 유한 관측 경로 & 참/거짓/미결정, 최초 확정 & 세 값과 최초 위반은 선행 기술; 본 논문의 증거 미상과 의미가 같지는 않음 \\
UPPAAL \cite{behrmann2004uppaal} & 오토마타와 상태 속성 & 속성 판정, 진단 경로 & 조건별 의무 상태와 최초 문제 기록을 실행하는 도구 \\
Scenic/VerifAI \cite{fremont2019scenic,dreossi2019verifai} & 상황 명세와 시험 시스템 & 생성 상황, 분석 결과 & 본 논문은 환경ㆍ제어 탐색 대신 고정 주석을 재생 \\
본 논문 & CoC 구절, 의무ㆍ조건 ID, 작성된 증거 & 모순/근거 부족/양립, 최초 사건ㆍ의무ㆍ사유 & 원문과 상태 판정의 연결을 실증; 자동 NLPㆍ자연 정확도 미검증 \\
\bottomrule
\end{tabularx}
\end{table*}

기술적 초점은 원문 구절과 의무ㆍ해제 조건ㆍ증거ㆍ문제 기록을 이어 검토 가능하게 만드는 데 있다. 기존 형식 검증 기술을 새로 발명했다고 주장하지 않으며, 다른 방법이 이러한 연결을 할 수 없다고 단정하지 않는다. 단순 FSM의 성공은 이 범위의 규칙을 간단한 상태 기계로도 실행할 수 있음을 보여 준다. UPPAAL은 표현한 모델의 상태 속성을 질의하고 진단 경로를 보존하는 역할을 맡는다.
\end{revision}

\section{데이터와 증거 범위}
\label{sec:data}

\begin{revision}

\subsection{자료 출처와 복원 범위}
\label{sec:provenance}
자연 자료는 nuScenes 기록에서 변환한 영상ㆍ자차 움직임과 사후 자동 생성한 CoC-Nusc 주석이다 \cite{caesar2020nuscenes,nvidia2026autolabeler}. 생성 과정은 기록 궤적에서 행동 종류를 계산하고 영상ㆍ행동 정보를 이용해 설명을 만든다. 따라서 CoC와 궤적의 일치는 독립적인 원인 정답이 아니다. 보관 목록에는 743장면ㆍ3,218사건ㆍ2,475인접쌍이 있으나 본문의 궤적 연결 분석은 94장면ㆍ403사건을 사용한다. 그중 CoC가 둘 이상인 90장면에서 309인접쌍을 구성했다. 사람 검토는 이 전체와 별도로 선정한 27장면 묶음이다.

CoC-Nusc의 배포 버전과 원 생성 실행별 checkpoint revision은 복원되지 않았다. 상류 README에는 모델 계열 Qwen/Qwen3.5-35B-A3B가 명시되어 있으나 개별 입력이 그 정확한 모델 revision으로 생성되었는지는 확인할 수 없다. 원 생성 prompt/config와 품질 점수 계산ㆍ상류 임계값도 미복원이다. 하류의 저장된 \texttt{quality\_pass} 불리언을 사용하는 선정 절차는 복원했다. 이번 보존 파일의 해시는 현재 입력 스냅샷을 식별하며 원 생성 이력을 새로 확정하지 않는다.

\begin{table}[t]
\centering
\color{revisionblue}
\bitablecaption{자료 출처의 확인 사실과 미복원 정보.}{Recovered provenance and information that remains unrecorded.}
\label{tab:provenance}
\footnotesize
\begin{tabularx}{\columnwidth}{L{.34\columnwidth}Y}
\toprule
항목 & 확인 범위 \\
\midrule
자연 입력 & 보관 파일ㆍ장면/사건 매핑ㆍ현재 해시 \\
생성 모델 & 상류 모델 계열 확인; 실행별 revision 미복원 \\
생성 prompt/config & 미복원; 새 prompt로 대신 채우지 않음 \\
품질 선정 & 저장된 quality\_pass 사용; 점수 계산/임계값 미복원 \\
작성 시나리오 & 원문ㆍ프레임 동시 작성, 기대 판정 실행 전 고정 \\
검사 실행 & 코드ㆍXMLㆍqueryㆍstdout/stderr 및 해시 보존 \\
\bottomrule
\end{tabularx}
\end{table}

보조 반응 분석의 선정은 $403\rightarrow290\rightarrow134$다. 품질 통과 290건 중 복합 행동 또는 방향이 불명확한 156건을 제외했다. 나머지 134건은 속도 변화 탐지 대상 130건과 사건 시점에 이미 정지ㆍ양보 조건을 만족한 4건이다. 기본 설정의 판정은 속도 반응 확인 125건, 이미 만족 4건, 제한 시간 내 미확인 5건이다. 이 선정은 309인접쌍의 문장 목록에 적용한 필터가 아니다.

\subsection{자연 검토 표본과 판정 기록}
\label{sec:selection}
후보 선택기는 이웃한 첫 문장에 \texttt{stop/yield} 계열, 둘째 문장에 \texttt{accelerate/proceed} 계열이 있는 쌍을 찾는다. 23쌍ㆍ15장면에서 장면별 가장 이른 창 하나를 선정했다. 비교군은 후보 장면을 제외하고, 첫 문장에 정지ㆍ양보 표현이 있으나 둘째 문장에는 가속ㆍ진행 표현이 없는 창이다. 장면당 가장 이른 창을 남기고, 후보와의 사건 수 차이, 정지/양보 종류 차이, 원 사건 위치 차이 순으로 대응시켰다. 동률은 seed 20260806과 창 ID의 해시로 결정했다. 비교 창을 중복 사용하지 않아 12장면이 선택되었다. 이 방식은 비무작위이며 특정 행동 전환을 더 많이 포함한다. 재분석에서 기존 27개 검토 ID와의 대응을 확인했다.

네 검토자의 원 CSV, 지침과 최종 합의 CSV를 사용했다. 지침은 시간 요구의 명시성, 요구 행동열과 문장 일관성을 구분하고, 외부 객체 상태가 필요하면 미상으로 남기도록 한다. 검토자의 전문성ㆍ교육 이력과 독립 중재자 참여는 충분한 기록이 없어 확인할 수 없다. 합의 CSV에는 27장면 모두 텍스트 모순 없음으로 저장되어 있지만 사유 메모는 0/27이다. 당시 논의 절차와 판단 변경의 원인을 재현 가능한 수준으로 복원하지 못했으며, 이번 작업에서 사람 재평가를 실시하지 않았다.

\subsection{작성 시나리오와 재현 자료}
\label{sec:artifacts}
통제 자료는 보행자, 자전거, 선행 차량, 교차 교통, 버스, 장애물의 6상황과 8변형, 중간 사건 0/2/6개의 모든 조합 144건이다. 실제 자연 문장에서 오류를 수집하거나 주입한 자료가 아니라 설계자가 영어 원문과 조건별 프레임을 함께 작성했다. 원문 구절의 시작ㆍ끝 위치, 사건ㆍ조건ㆍ의무 ID와 기대 판정이 저장되어 있다. 사례와 구현에 같은 설계자가 관여했으므로 독립 외부 평가가 아니다.

재현 자료는 프로젝트 소유 산출물 \texttt{kiee-review-revision-001} 아래에 보존했다. 공개 분류의 \texttt{controlled-scenarios-2026-10-01-001}에는 입력ㆍ최초 결과ㆍ48개 실패ㆍ모델 로그, \texttt{terminal-evidence-2026-10-01-001}에는 종료 질의 진단, \texttt{condition-model-2026-10-01-001}에는 동일 사례 재시험이 있다. 사람 원 CSVㆍ원문 및 기존 자료는 restricted 분류를 유지한다. 파일을 보존했다는 사실과 외부 공개를 완료했다는 사실은 구분한다. 이번 개정은 외부 공개를 수행하지 않았다.
\end{revision}

\section{문제 정의}
\label{sec:problem}

\begin{revision}

장면의 사건을 $E_s=(e_0,\ldots,e_{n-1})$로 두고 자료상 연속성이 확인된 구간에서만 의무를 이어받는다. 사건은 행동, 새 의무, 조건별 증거와 원문 구절을 가진다. 의무 $o$는 ID, 대상과 해제 조건 ID $c(o)$에 연결된다. 조건 증거 $q_i(c)\in\{\mathsf{T},\mathsf{F},?\}$에서 참은 해당 조건 충족, 거짓은 미충족 확인, $?$는 알 수 없음을 뜻한다. 새 행동 근거 B는 ID가 다른 A의 해제 증거가 아니다.

\subsection{주석 문제와 검사기 처리 오류의 분리}
\label{sec:error-types}
P1--P3은 입력의 행동ㆍ요구 관계에 관한 문제이고 P4는 검사기의 상태 처리 문제다. P5는 별도의 속도 반응 시간 검사다. 기존 40쌍을 이 구분으로 재집계한다.

\begin{align}
C_1(i)&=\mathrm{GO}_i\land\exists o\in A_i^-:q_i^*(c(o))=\mathsf{F},\label{eq:hold-conflict}\\
C_2(i)&=\exists o:\mathrm{release}_i(o)\land q_i(c(o))=\mathsf{F},\\
C_3(i)&=\mathrm{GO}_i\land p_i^*=\mathsf{F},\\
F_4(i)&=\exists o\in A_{i-1}:q_i^*(c(o))=\mathsf{T}\land o\in A_i,\\
M(i)&=\mathrm{GO}_i\land\exists o\in A_i^-:q_i^*(c(o))=? .\label{eq:missing-evidence}
\end{align}

$A_i^-$는 새 의무와 해제 증거를 반영한 뒤 현재 행동 직전의 활성 의무 집합이다. $q_i^*$는 유효한 최근 증거를, $p_i^*$는 저장된 선행조건 증거를 뜻한다. P1은 미해제가 확인된 의무와 진행의 충돌, P2는 미충족 조건을 충족했다고 처리한 조기 해제, P3은 선행조건 미충족 상태의 진행이다. $F_4$는 해제가 확인되었는데도 해당 의무를 상태에 남기는 처리 오류다. $M$은 진행 시 해제 근거 부족이며 확정 모순에 포함하지 않는다.

기존 검사기의 P2 출력은 진행 사건의 \texttt{release=false}에서, P4 출력은 같은 대상의 정지ㆍ양보 반복과 \texttt{release=true}가 함께 주어진 시험에서 발생한다. 이 구현상 표시는 실제 사람이 조기 해제했다고 독립 확인한 결과가 아니다. P4 변형 10건은 의도적으로 만든 처리 시험으로 보고한다. 보완 모델은 별도의 해제 명령이나 P4 고장 주입을 받지 않고 조건 ID의 참 증거로 정상 상태를 갱신한다. 따라서 보완 모델의 144/144를 P1--P4 전체의 새 성능이라고 서술하지 않는다.

\subsection{조건별 갱신과 판정 의미}
\label{sec:semantics}
보완 모델의 갱신은 다음과 같다. 미상 입력은 최근의 알려진 증거를 덮지 않는다.
\begin{equation}
q_i^*(c)=\begin{cases}q_i(c),&q_i(c)\ne ?,\\q_{i-1}^*(c),&q_i(c)=?,\end{cases}
\end{equation}
\begin{equation}
A_i^-=(A_{i-1}\cup O_i)\setminus\{o:q_i^*(c(o))=\mathsf{T}\},\qquad A_i=A_i^- .
\label{eq:active-update}
\end{equation}
이 보존 규칙은 같은 대상ㆍ조건이 유지되고 증거의 유효 구간이 계속된다는 작성 자료의 가정에 의존한다. 실제 객체의 추적ㆍ증거 만료ㆍ재등장은 모델링하지 않는다. 한 번 해제된 의무는 같은 ID의 재생성 입력이 없으면 다시 활성화하지 않는다.

정지나 대기 중 해제 여부가 미상이라는 이유만으로 문제를 기록하지 않는다. 진행할 때 남은 의무가 거짓이면 모순, 미상이면 근거 부족이다. 복수 의무는 각각의 조건을 확인하며 하나의 해제가 다른 의무를 지우지 않는다. 같은 사건에 모순과 근거 부족이 함께 있으면 모순을 최종 판정으로 우선하되 각각의 기록을 보존한다. 전체 판정은 누적 모순이 있으면 모순, 그렇지 않고 진행 중 근거 부족이 있었으면 미상, 둘 다 없으면 표현한 의무와 양립이다. 마지막 판정은 모든 위험이 없다는 뜻이 아니다.

장면 끝에서 미래의 해제를 보지 못했다는 사실만으로 모순을 추가하지 않는다. 미상 상태에서 진행하지 않은 대기열은 현재 검사 범위에서 문제 없음으로 끝날 수 있다. 근거 없이 다음 장면과 연결하지 않으며 다음 장면은 초기화한다. 임의의 미래 완료 의무나 제한 시간을 추가한 liveness 명세는 이번 중심 검사에 없다.
\end{revision}

\section{\rev{제안 방법과 구현}}
\label{sec:method}

\begin{revision}

\subsection{자동 텍스트 변환과 작성 프레임}
\label{sec:extraction}
기존 자동 변환은 LLM 호출이나 수작업 정답 입력을 쓰지 않는 정규식 기반 영어 파서다. \texttt{stop/halt/hold}, \texttt{yield/decelerate/slow down}, \texttt{accelerate/proceed/resume/move forward}, \texttt{maintain/keep speed}를 행동 종류로 묶는다. 명령형ㆍ조동사ㆍ주어 문맥을 확인하여 단순 언급을 행동으로 해석하는 것을 제한한다. \texttt{then/and then}으로 단계 순서를 나누되 \texttt{until then}은 단계 경계로 자르지 않는다. \texttt{until} 뒤 구절은 해제 조건, \texttt{once/after} 뒤 구절은 계기로 보존하고 조건 구절 내부의 행동 단어는 제외한다.

정지ㆍ양보 대상이 없으면 \texttt{UNKNOWN\_MISSING\_TARGET}, 명확한 순서 없는 복수 행동이면 \texttt{UNKNOWN\_AMBIGUOUS\_ORDER}로 둔다. 명시적 단계가 각각 단일 행동일 때만 다음 행동을 연결한다. 부정 표현을 참ㆍ거짓 관측 증거로 읽거나 동의 표현을 범용 의미로 해석하는 기능은 없다. \texttt{release\_known}과 \texttt{satisfaction\_known}은 모두 false로 시작한다. 문구를 찾은 것과 해당 조건이 실제로 충족되었음을 안 것은 다르다.

작성 프레임은 설계자가 원문과 동시에 만든다. \texttt{action}, \texttt{obligation\{id,target,condition\_id\}}, \texttt{evidence\{condition\_id,value\}}, \texttt{precondition}과 \texttt{source\_spans}를 가진다. True/False는 작성 문장이 명시한 조건 상태이며 사람의 독립 센서 정답이 아니다. 기존 모델에 입력할 때 조건 A를 단일 해제 채널에 투영하고 B/OTHER를 A로 바꾸지 않는다. 두 번째 의무 C의 별도 해제 정보를 담지 못한 손실도 기록한다. 조건별 모델은 프레임의 각 의무를 직접 인코딩한다.

\begin{table*}[t]
\centering
\color{revisionblue}
\bitablecaption{가상 보행자 사례의 원문--작성 프레임--모델 상태--판정 연결(중간 사건 0개).}{End-to-end authored pedestrian example with zero intervening events.}
\label{tab:worked-example}
\scriptsize
\begin{tabularx}{\textwidth}{L{.20\textwidth}L{.32\textwidth}L{.25\textwidth}Y}
\toprule
사건/변형 & 실제 작성 원문 구절 & 프레임과 상태 & 출력 \\
\midrule
공통 e0 & Stop until the pedestrian has passed. The required mirror check has been completed. & STOP; 의무 A, 대상 pedestrian; 해제조건 A; precondition=T; active[A]=true, known[A]=? & 의무 시작; 문제 없음 \\
미해제 e1 & The pedestrian is still crossing. The traffic signal is green. Proceed. & GO; evidence A=F, B=T; active[A]=true & 모순; e1, A, UNRELEASED\_OBLIGATION \\
증거 누락 e1 & The traffic signal is green. Proceed. & GO; evidence B=T; known[A]=? & 근거 부족; e1, A, RELEASE\_EVIDENCE\_MISSING \\
정상 해제 e1 & The pedestrian has passed. The traffic signal is green. Proceed. & GO; evidence A=T, B=T; active[A]=false & 해당 의무와 양립 \\
\bottomrule
\end{tabularx}
\end{table*}

표~\ref{tab:worked-example}의 원문은 저장된 작성 사례에서 가져왔다. 미해제 e1의 \texttt{Proceed}는 문자 위치 [63,70), A 증거는 [0,33), B 증거는 [34,62)에 연결된다. 사건 ID는 \texttt{pedestrian-held\_local\_conflict-gap0:e1}이다. 누락ㆍ정상 변형도 각각의 case ID에 e1을 연결한다. 실제 입력은 UTF-8 구절의 문자열 인덱스이며 이 예의 원문은 ASCII다. 자동 파서에 같은 문장을 넣는 경로는 증거를 확정하지 않으므로 작성 프레임의 성공과 혼합하지 않는다.

\subsection{기존 검사기와 보완 모델의 차이}
\label{sec:models}
기존 Python은 활성 의무 한 개, 행동 4종, 상태 6종과 누적 오류ㆍ미상 사유를 관리한다. 단계마다 해석 미상, 알려진 증거 갱신, 오류 규칙, 상태 전이, 누적 출력 순으로 처리한다. 복수 의무는 \texttt{OVERLAPPING\_OBLIGATION}, 대기 중 해제 미상도 누적 미상으로 남을 수 있다. 기존 UPPAAL은 같은 전이 규약과 해시를 사용한다. 이 두 경로는 동결하여 최초 실패를 보존했다.

보완 모델은 실패 분석 뒤 별도 개발했다. 각 의무의 \texttt{active[K]}, 조건별 \texttt{known[K]}(-1/0/1), 선행조건, 사건 인덱스, \texttt{bad/missing} 누적 플래그와 최초 위치ㆍ의무를 둔다. 입력의 의무 등록 순서가 배열 순서이며 선행조건 위반을 먼저, 의무를 등록 순서대로 판정한다. 같은 사건의 여러 위반을 개별 사유로 누적하되 최초 의무 출력은 이 우선순위로 하나를 선택한다. 모든 문제의 전체 순위나 모든 관련 의무 목록을 반환하는 기능은 아니다.

\begin{table}[t]
\centering
\color{revisionblue}
\bitablecaption{보완 모델의 조건별 상태 갱신과 기록 순서.}{Update and recording order in the condition-indexed model.}
\label{tab:transition}
\scriptsize
\begin{tabularx}{\columnwidth}{L{.27\columnwidth}Y}
\toprule
단계 & 실행과 기록 \\
\midrule
1. 등록 & 현재 사건의 새 의무를 active에 넣음 \\
2. 증거 & 알려진 선행조건과 조건별 증거만 갱신; 미상은 이전 값 보존 \\
3. 해제 & 각 의무의 known=1이면 해당 active만 false \\
4. 진행 & GO일 때 precondition=0이면 bad; 이후 등록 순서로 active/known=0이면 bad, -1이면 missing \\
5. 최초 기록 & 최초 bad/missing의 사건ㆍ의무 인덱스를 한 번만 기록; 이후 플래그 누적 \\
6. 종료 & 사건 배열 끝에서 Done; 미래 해제의 부재만으로 오류를 추가하지 않음 \\
\bottomrule
\end{tabularx}
\end{table}

\begin{lstlisting}[caption={보완 모델의 핵심 실행 순서(실제 XML의 step 함수 요약).},label={lst:condition-step}]
activate(creates[idx]);
update_known_precondition_if_present();
for j in obligation_registration_order:
    if observation(idx,j) >= 0: known[j] = observation(idx,j)
    if active[j] and known[j] == 1: active[j] = false
if actions[idx] == GO:
    record_bad_if(precondition == 0, PRECONDITION)
    for j in obligation_registration_order:
        record_bad_if(active[j] and known[j] == 0, j)
        record_missing_if(active[j] and known[j] == -1, j)
idx = idx + 1
\end{lstlisting}

WAIT는 작성 자료에서 유지ㆍ대기 단계를 나타내는 기호다. 실제 속도나 이동을 적분하지 않으므로 \texttt{Maintain speed}라는 문구를 물리적으로 안전한 정지라고 보장하지 않는다. 이 단계는 명시된 GO 전환 위반을 검사하기 위한 중간 사건이다.

\subsection{UPPAAL 상태ㆍ질의와 추적}
\label{sec:queries}
보완 모델은 committed 위치 Run에서 \texttt{idx<N}일 때 \texttt{step(),idx++}로 자기 전이하고, \texttt{idx==N}이면 Done으로 이동한다. committed 처리로 한 사건의 갱신을 끝내고 다음 사건을 읽는다. 정량 시간 시계는 없다. 다음 다섯 질의를 실제 verifyta 5.0.0에 실행했다. 경로ㆍ상태 양화사의 의미는 공식 문서에 따른다 \cite{uppaalqueries}.

\begin{lstlisting}[caption={보완 모델의 실제 queries.q.},label={lst:queries}]
A[] not bad
A[] not missing
A<> Observer.Done
E<> Observer.Done && bad_hold
E<> Observer.Done && bad_pre
\end{lstlisting}

첫 질의의 실패는 모순, 둘째 질의만 실패하면 근거 부족이다. 셋째는 전체 배열의 처리 완료를 확인하고 넷째ㆍ다섯째는 종료 상태에서 미해제 의무 및 선행조건 위반 사유를 확인한다. 진단 경로의 \texttt{first\_bad\_idx/first\_missing\_idx}, 관련 의무 인덱스를 원 사건/의무 ID로 되돌린다. 조건 B는 A에 대응하는 관측 배열에 들어가지 않는다.

기존 모델은 P1--P4 플래그별 \texttt{A[] not flag}, \texttt{A[] not sticky\_unknown}, \texttt{A<> Observer.Done}을 사용했다. 최초 반례 접두 경로가 뒤에 생긴 미상 사유를 담지 못한 24건은 종료 상태 추가 질의로 진단했다. 원 파서는 변경하지 않았다. 두 기존 구현의 일치는 고정 규칙의 구현 일치성이지 그 규칙의 독립 타당성이나 자연어 해석의 정답성이 아니다.

보조 반응 모델은 기록 사건과 자차 속도 반응을 재생하고 제한 시간ㆍ반복 시계ㆍ출처 덮어쓰기를 검사한다. 중심 조건별 모델과 질문 및 입력이 다르므로 그 성공을 중심 문장 검사의 안전 검증으로 합치지 않는다.
\end{revision}

\section{\rev{평가 설계}}
\label{sec:experiments}

\begin{revision}

\subsection{기존 회귀시험과 자연 검토}
\label{sec:legacy-test}
기존 40개의 기준--변형 쌍은 각 규칙에 10쌍씩 배정하고 기대 판정이 나오는 쌍만 남긴 규칙 실행 회귀시험이다. P1--P3의 주석/구조화 사건 시험 30쌍과 P4 처리 시험 10쌍을 분리한다. 이 회귀 결과를 자연 오류 precision/recall/F1로 제시하지 않는다. 이전 요구를 제거한 현재 사건 검사는 기억 효과를 확인하는 ablation이다.

자연 검토는 선정된 27장면에서 네 검토자의 시간 요구ㆍ행동열ㆍ텍스트 일관성 범주를 재집계한다. 원 라벨은 변경하지 않고 Fleiss' $\kappa$와 평균 쌍별 일치율, 검토자별 분포를 보고한다. 다수가 선택한 라벨이나 합의 CSV를 새 외부 정답으로 삼지 않는다. 이번 개정에서 사람 재평가ㆍ독립 중재와 자연 오류 정답 수집을 추가하지 않았다.

\subsection{144개 통제 시나리오의 최초 실행}
\label{sec:scenario-design}
6상황의 8변형을 중간 사건 0/2/6개와 교차시켜 144개를 모두 실행했다. 각 변형은 18건이고 기대 분포는 모순 72, 양립 36, 근거 부족 36이다. 기대 판정ㆍ최초 사건ㆍ의무ㆍ사유는 최초 실행 전에 고정했으며 결과에 따른 선별은 하지 않았다. 원문과 프레임은 동시 작성했으므로 자연 문장 자동 추출 시험과 구분한다. 첫 결과에 맞지 않은 48개도 모두 보존했다.

\begin{table}[t]
\centering
\color{revisionblue}
\bitablecaption{시나리오 변형과 고정된 기대 판정(변형별 18건).}{Scenario variants and fixed expected verdicts, 18 cases per variant.}
\label{tab:scenario-variants}
\scriptsize
\begin{tabularx}{\columnwidth}{L{.45\columnwidth}Y}
\toprule
변형 & 기대 판정의 근거 \\
\midrule
현재 미해제 후 진행 & A=F, GO: 모순 \\
이전 구간 미해제 후 진행 & A=F의 유효 구간 유지, GO: 모순 \\
정상 해제 & A=T 후 GO: 양립 \\
해제 증거 누락 & A=?, B=T, GO: 근거 부족 \\
다른 대상 해제 & OTHER=T, A=?, GO: 근거 부족 \\
행동 순서 위반 & A=T이나 선행 확인=F, GO: 모순 \\
복수 의무 일부 해제 & A=T, C=F, GO: 모순 \\
복수 의무 모두 해제 & A=T, C=T, GO: 양립 \\
\bottomrule
\end{tabularx}
\end{table}

비교 경로는 동결된 기존 Python 전체 이력 검사, 현재 사건 검사, 이전 1/3/5개 사건 창, 별도 FSM, 기존 UPPAAL, 기존 자동 텍스트 경로다. 창 비교는 기존 검사기에 잘린 이력을 제공한 ablation이며 독립 경쟁 알고리즘은 아니다. 구조화 기존 경로에는 단일 의무 투영, FSM에는 조건별 프레임을 제공한다. 따라서 그 차이에는 기억 길이뿐 아니라 의무 표현 용량과 미상 누적 정책도 포함된다. 단일 의무를 지원하는 경로의 복수 의무 도전 사례도 분모에서 제외하지 않는다.

FSM은 active 의무 집합과 조건별 알려진 증거를 갱신하는 별도 단순 구현이다. 작성된 조건 상태를 입력받으므로 자연어 의미를 독립 판정하는 oracle이 아니다. 기존 자동 경로에는 원문만 제공하고 증거 정답을 넣지 않았다. 세 분류 일치와 모순 72건을 분모로 한 최초 위반 위치 일치를 보고하며, 모순 사례를 미상으로 남긴 것도 실패 분모에 포함한다.

\subsection{동일 사례 재시험과 구현 일치성}
\label{sec:retest-design}
48개 최초 불일치의 원인을 분석한 뒤 조건별 UPPAAL 모델을 개발했다. 기존 코드와 기대 판정은 바꾸지 않고 같은 \texttt{cases.json}을 재사용했다. 개발 후 동일 사례 재시험이며 개발 미사용 독립 시험이 아니다. 모순 72건의 최초 위반과 모순/근거 부족 108건의 문제 위치ㆍ관련 의무ㆍ사유를 비교했다. 성공만 보고하지 않도록 최초 시험과 재시험을 별도 표에 둔다.

기존 Python--UPPAAL 비교는 판정ㆍ오류 종류ㆍ최초 위치ㆍ미상 사유를 항목별로 대조한다. 사유가 다른 사례는 최초 반례와 종료 상태를 따로 읽어 모델 의미 문제와 출력 범위 문제를 구분한다. 전체 상태 경로의 동등성이나 규칙 자체의 독립 타당성을 검증한 시험은 아니다.

\subsection{보조 분석과 통계의 범위}
\label{sec:aux-design}
속도 반응 분석은 $W=0.5$초 구간에서 속도 변화량 $\delta=0.1$~m/s 이상이고 방향 일관 비율 $\rho=0.7$ 이상인 시작점을 찾으며 $D=3$초 제한과 비교한다. 정지ㆍ양보의 $v\leq0.5$~m/s 이미 만족 규칙은 초기 후보 검토 뒤 추가되었다. 이 값들은 법규ㆍ차량 안전 기준에서 도출한 것이 아니라 탐색 설정이다. 기존 27설정의 130건을 항상 검출ㆍ설정 의존ㆍ항상 미검출로 재집계했다. 반복 CoC의 최초 시각 보존과 재시작, 장면 경계 연결 처리도 별도 보조 결과로 유지한다.

144건은 동일 문장 틀과 변형을 교차한 결정적 조합 시험이다. 자연 모집단의 무작위 표본이나 독립 시행이 아니므로 모집단 오류율과 95\% 신뢰구간을 산출하지 않는다. 자연 검토도 확정된 독립 오류 정답이 없어 TP/FP/TN/FN, sensitivity, specificity, precision/recall/F1을 계산하지 않는다. 미상을 정상으로 합치지 않고 실제 분류 건수를 함께 보고한다.
\end{revision}

\section{\rev{결과}}
\label{sec:results}

\begin{revision}

\subsection{기존 40쌍의 규칙 실행 결과}
\label{sec:regression-results}
표~\ref{tab:regression}은 자연 오류 성능이 아닌 회귀 결과다. 기존 전체 이력 검사는 P1--P3 30/30과 P4 처리 시험 10/10, 현재 사건 검사는 각각 20/30과 0/10이었다. 기준 사례의 확정 모순 표시는 양쪽 모두 0/40이다. 합계 40/40과 20/40을 일반 오류 검출 정확도로 해석하지 않는다.

\begin{table}[t]
\centering
\color{revisionblue}
\bitablecaption{기존 규칙 실행 회귀시험: 주석 문제와 처리 오류를 분리.}{Legacy rule-execution regression, separating annotation-related rules and processing faults.}
\label{tab:regression}
\footnotesize
\begin{tabularx}{\columnwidth}{Ycc}
\toprule
시험 & 전체 이력 & 현재 사건 \\
\midrule
P1 미해제 의무 충돌 & 10/10 & 0/10 \\
P2 조기 해제 & 10/10 & 10/10 \\
P3 행동 순서 & 10/10 & 10/10 \\
P4 검사기 처리 오류 & 10/10 & 0/10 \\
기준 사례 모순 표시 & 0/40 & 0/40 \\
\bottomrule
\end{tabularx}
\end{table}

\subsection{최초 144개 작성 시나리오}
\label{sec:first-results}
표~\ref{tab:first-results}의 전체 이력 경로는 작성 프레임을 기존 단일 의무 형식으로 투영한 결과이며 96/144가 기대 판정과 일치했다. 확정 모순 72개 중 최초 위반 일치는 54/72다. 현재 사건은 72/144, 이전 1/3/5개는 84/90/90개가 일치했다. FSM은 144/144, 최초 위반 72/72였다. UPPAAL이 단순 비교기보다 정확하다는 근거는 없다.

\begin{table*}[t]
\centering
\color{revisionblue}
\bitablecaption{고정된 144개 작성 시나리오의 최초 결과. 간격별 분모는 각각 48; 조건별 보완 재시험은 이 표에 포함하지 않는다.}{First-run results on all 144 authored cases, with 48 cases per gap. The condition-indexed retest is excluded.}
\label{tab:first-results}
\scriptsize
\begin{tabularx}{\textwidth}{Yccccc}
\toprule
경로 (입력) & 세 분류 일치 & 최초 위반 & 간격0 & 간격2 & 간격6 \\
\midrule
기존 전체 이력 (작성 프레임 투영) & 96/144 & 54/72 & 36/48 & 30/48 & 30/48 \\
현재 사건 (투영) & 72/144 & 36/72 & 24/48 & 24/48 & 24/48 \\
이전 1개 (투영) & 84/144 & 42/72 & 36/48 & 24/48 & 24/48 \\
이전 3개 (투영) & 90/144 & 48/72 & 36/48 & 30/48 & 24/48 \\
이전 5개 (투영) & 90/144 & 48/72 & 36/48 & 30/48 & 24/48 \\
별도 FSM (조건별 프레임) & 144/144 & 72/72 & 48/48 & 48/48 & 48/48 \\
기존 자동 텍스트 & 36/144 & 0/72 & 12/48 & 12/48 & 12/48 \\
기존 UPPAAL (투영) & 96/144 & 54/72 & 36/48 & 30/48 & 30/48 \\
\bottomrule
\end{tabularx}
\end{table*}


기존 구조화 전체 이력 경로의 혼동 분포는 기대 모순 72건 중 모순 54ㆍ미상 18, 기대 양립 36건 중 양립 6ㆍ미상 30, 기대 근거 부족 36건 중 미상 36이다. 따라서 정상 오탐 0/36만 보고하면 정상 대부분을 판정하지 못한 사실을 숨긴다. 기존 자동 텍스트 경로는 모든 144건을 미상으로 판정했다. 기대 미상 36건만 일치하며 구조화 검사 성공을 자동 추출 성공으로 바꿀 수 없다.

변형별로 기존 경로는 현재 미해제ㆍ이전 미해제ㆍ증거 누락ㆍ다른 대상 해제ㆍ행동 순서 위반에서 각각 18/18이었다. 정상 해제는 6/18이고 복수 일부 해제와 모두 해제는 각각 0/18이다. 간격별 일치 건수와 최초 위반 위치는 표~\ref{tab:first-results}에 함께 제시했다. 창을 늘려도 단일 의무 표현과 누적 미상 처리의 한계는 남는다.

\subsection{실패 분석과 개발 후 동일 사례 재시험}
\label{sec:retest-results}
최초 불일치 48건은 복수 일부 해제 18, 복수 모두 해제 18, 중간 대기 후 정상 해제 12다. 복수 사례에서는 단일 의무 투영과 중첩 처리로 별도 조건을 표현하지 못했다. 정상 해제 사례에서는 대기 중 미상이 누적되어 이후 명시적인 정상 해제로도 최종 미상이 남았다. 이 실패는 자연 오류가 아니라 작성 의미와 기존 표현/정책의 불일치다.

\begin{table}[t]
\centering
\color{revisionblue}
\bitablecaption{48개 최초 불일치와 보완의 대상.}{All 48 initial mismatches and the corresponding refinement.}
\label{tab:failures}
\scriptsize
\begin{tabularx}{\columnwidth}{L{.36\columnwidth}cY}
\toprule
유형 & 건수 & 보완 내용 \\
\midrule
복수 일부 해제 & 18 & C의 조건/의무를 A와 별도로 유지 \\
복수 모두 해제 & 18 & 각 조건의 참 증거로 각각 해제 \\
중간 대기 후 정상 해제 & 12 & 미상 대기와 근거 없이 진행한 사건 구분 \\
\bottomrule
\end{tabularx}
\end{table}

\begin{table}[t]
\centering
\color{revisionblue}
\bitablecaption{개발 후 동일 사례 재시험: 조건별 UPPAAL 모델(독립 시험 아님).}{Post-development same-case retest of the condition-indexed UPPAAL model, not an independent test.}
\label{tab:retest}
\footnotesize
\begin{tabularx}{\columnwidth}{Yc}
\toprule
대조 항목 & 기대 기록과 일치 \\
\midrule
최종 세 분류 & 144/144 \\
모순의 최초 위반 사건 & 72/72 \\
문제 위치(모순+근거 부족) & 108/108 \\
관련 의무 & 108/108 \\
문제 사유 & 108/108 \\
\bottomrule
\end{tabularx}
\end{table}

보완 결과는 조건별 표현과 문제 추적을 실행할 수 있음을 보여 준다. 기대 판정을 바꾸지는 않았으나 모델 개발에 같은 실패 사례를 사용했다. 따라서 독립 일반화 성능이 아니며 R3의 개발 미사용 시험 요구를 충족한 결과로 제시하지 않는다. 새 프레임 모델이 자동 NLP 추출을 개선한 것도 아니다.

\subsection{기존 Python--UPPAAL 출력 대조}
\label{sec:agreement-results}
표~\ref{tab:implementation}에서 판정ㆍ오류 종류ㆍ최초 위치는 기존 두 구현의 144건 전체에서 일치했다. 미상 사유는 24건이 달랐다. 최초 위반 반례의 접두 경로가 뒤에 누적된 미상 사유를 담지 못한 경우였다. 종료 상태의 추가 질의로 뒤의 사유까지 확인하면 144/144가 일치했다. 원 파서와 최초 로그는 보존했다. 이 출력 범위 진단을 48개 의미 불일치의 해결과 혼합하지 않는다.

\begin{table}[t]
\centering
\color{revisionblue}
\bitablecaption{기존 Python--UPPAAL 구현 일치성 및 종료 질의 진단.}{Legacy Python--UPPAAL output agreement and terminal-query diagnosis.}
\label{tab:implementation}
\footnotesize
\begin{tabularx}{\columnwidth}{Yc}
\toprule
대조 항목 & 일치 사례 \\
\midrule
최종 판정 / 오류 종류 / 최초 위치 & 각 144/144 \\
원 반례 접두 경로의 미상 사유 & 120/144 \\
종료 상태 추가 질의의 미상 사유 & 144/144 \\
\bottomrule
\end{tabularx}
\end{table}

\subsection{자연 검토의 일치도와 기록 한계}
\label{sec:human-results}
\begin{table}[t]
\centering
\color{revisionblue}
\bitablecaption{네 검토자 원 판정의 재집계(27장면).}{Reanalysis of four reviewers' original labels on 27 scene groups.}
\label{tab:agreement}
\footnotesize
\begin{tabularx}{\columnwidth}{Yccc}
\toprule
항목 & $\kappa$ & 쌍별 일치율 & 불일치 장면 \\
\midrule
시간 요구 & 0.202892 & 0.697531 & 14/27 \\
행동열 & 0.291391 & 0.339506 & 23/27 \\
텍스트 일관성 & 0.027301 & 0.407407 & 23/27 \\
\bottomrule
\end{tabularx}
\end{table}
\begin{table}[t]
\centering
\color{revisionblue}
\bitablecaption{검토자별 텍스트ㆍ시간 요구 판정 분포(각 27건, 익명 R1--R4).}{Per-reviewer textual and temporal-requirement label distributions, 27 ratings per field.}
\label{tab:reviewer-distribution}
\footnotesize
\begin{tabularx}{\columnwidth}{Ycccc}
\toprule
판정 & R1 & R2 & R3 & R4 \\
\midrule
텍스트 모순 없음 & 9 & 18 & 10 & 25 \\
텍스트 모순 & 10 & 0 & 1 & 1 \\
외부 근거 필요 & 4 & 4 & 5 & 0 \\
문장 모호 & 1 & 0 & 0 & 0 \\
해당 없음 & 3 & 5 & 11 & 1 \\
\midrule
시간 요구 있음 & 18 & 22 & 15 & 26 \\
시간 요구 없음 & 9 & 5 & 11 & 1 \\
시간 요구 모호 & 0 & 0 & 1 & 0 \\
행동열의 서로 다른 라벨 수 & 20 & 15 & 18 & 10 \\
\bottomrule
\end{tabularx}
\end{table}

텍스트 $\kappa=0.027301$, 시간 요구 $\kappa=0.202892$, 행동열 $\kappa=0.291391$을 재현했다. 텍스트의 평균 쌍별 일치율은 0.407407이고 23/27장면에서 네 명의 판정이 달랐다. 표~\ref{tab:reviewer-distribution}에서 텍스트 모순 판정은 검토자별 10/0/1/1건, 해당 없음은 3/5/11/1건으로 크게 다르다. 이는 라벨 적용과 해석의 불안정성을 보여 주지만 원인을 전문성이나 교육 차이로 확정할 수는 없다.

원 메모에는 ``시간 요구가 없는 행동 조절'', ``해제 조건과 행동 순서의 구별'', ``선행 관계를 연결할 정보 부족'' 등의 해석이 남아 있다. 이는 저장된 개별 의견이지 합의 원인의 전수 분석이 아니다. 현재 재분석으로 불일치의 질적 원인별 건수나 합의 형성 절차를 새로 확정하지 않았다. 합의 CSV는 모순 0/27이나 사유 기록도 0/27이다. 이를 객관적인 자연 정답으로 쓰지 않고 재평가ㆍ독립 중재가 미실시임을 유지한다.


\begin{table*}[t]
\centering
\color{revisionblue}
\bitablecaption{원 판정의 주요 불일치 예시. 개별 메모는 합의 사유나 확정 원인 분류가 아니다.}{Examples of original disagreements; individual notes are not consensus rationales or adjudicated causes.}
\label{tab:disagreement-examples}
\scriptsize
\begin{tabularx}{\textwidth}{lL{.29\textwidth}Y}
\toprule
익명 장면 & 네 원 판정(R1/R2/R3/R4) & 보존된 개별 의견의 요지 \\
\midrule
G01 & 모순/해당 없음/해당 없음/모순 없음 & R3: 속도 조절 외 시간 요구가 없다고 해석 \\
G04 & 모순/외부 근거 필요/외부 근거 필요/모순 없음 & R3: 신호가 바뀐 뒤 출발했는지 알 수 없어 증거 필요 \\
G22 & 모순 없음/모순 없음/외부 근거 필요/모순 & R3: 보행자 통과 여부에 외부 근거 필요; R4: 정지ㆍ양보 요구로 해석 \\
\bottomrule
\end{tabularx}
\end{table*}

\subsection{궤적과 보조 반응 결과}
\label{sec:trajectory-results}
27장면의 궤적 비교는 일치 2, 불일치 7, 미상 18이며 재실행 기록과 장면별 판정 차이는 0이었다. 표~\ref{tab:trajectory-cases}는 각 장면의 저장 사유를 익명 순번으로 제시한다. 행동 일치가 있더라도 해제나 단계 시점 증거가 부족한 사건이 있으면 미상이 될 수 있다. 종방향 행동 불일치가 있는 7건은 문장 사이의 모순이나 위험 주행의 독립 정답이 아니다.

\begin{table*}[t]
\centering
\color{revisionblue}
\bitablecaption{27장면 궤적 판정의 사례별 사유. 익명 G01--G27은 보존 재분석 행 순서; 사유는 사건 건수이며 한 장면에 중복될 수 있다.}{Case-level reasons for trajectory comparison; anonymous IDs follow preserved reanalysis row order. Counts refer to events and can overlap within a scene.}
\label{tab:trajectory-cases}
\scriptsize
\begin{tabularx}{\textwidth}{lYrrrrrr}
\toprule
장면 & 장면 판정 & 행동 일치 & 종방향 불일치 & 해석 미상 & 행동 미지원 & 관측 미상 & 시점/해제 부족 \\
\midrule
G01 & 미상 & 0 & 0 & 0 & 4 & 2 & 0 \\
G02 & 미상 & 0 & 0 & 2 & 1 & 0 & 0 \\
G03 & 불일치 & 1 & 1 & 2 & 0 & 0 & 1 \\
G04 & 미상 & 1 & 0 & 1 & 1 & 0 & 1 \\
G05 & 불일치 & 0 & 1 & 1 & 4 & 0 & 0 \\
G06 & 미상 & 2 & 0 & 0 & 0 & 0 & 2 \\
G07 & 미상 & 3 & 0 & 1 & 0 & 0 & 0 \\
G08 & 미상 & 4 & 0 & 0 & 3 & 0 & 2 \\
G09 & 미상 & 3 & 0 & 0 & 1 & 0 & 2 \\
G10 & 불일치 & 0 & 1 & 0 & 1 & 0 & 0 \\
G11 & 미상 & 3 & 0 & 1 & 1 & 0 & 0 \\
G12 & 미상 & 3 & 0 & 1 & 0 & 0 & 2 \\
G13 & 불일치 & 1 & 1 & 1 & 0 & 0 & 1 \\
G14 & 미상 & 1 & 0 & 0 & 0 & 0 & 2 \\
G15 & 미상 & 0 & 0 & 0 & 4 & 0 & 1 \\
G16 & 일치 & 3 & 0 & 0 & 0 & 0 & 0 \\
G17 & 미상 & 0 & 0 & 0 & 2 & 0 & 1 \\
G18 & 일치 & 3 & 0 & 0 & 0 & 0 & 0 \\
G19 & 미상 & 0 & 0 & 0 & 2 & 0 & 1 \\
G20 & 미상 & 0 & 0 & 4 & 3 & 0 & 0 \\
G21 & 미상 & 0 & 0 & 0 & 2 & 0 & 0 \\
G22 & 미상 & 1 & 0 & 0 & 3 & 0 & 0 \\
G23 & 불일치 & 4 & 1 & 0 & 0 & 0 & 0 \\
G24 & 불일치 & 2 & 3 & 0 & 0 & 0 & 0 \\
G25 & 불일치 & 0 & 1 & 1 & 1 & 0 & 1 \\
G26 & 미상 & 0 & 0 & 0 & 6 & 0 & 0 \\
G27 & 미상 & 2 & 0 & 0 & 2 & 0 & 0 \\
\bottomrule
\end{tabularx}
\end{table*}


보조 반응 134건의 기본 판정은 $125+4+5$다. 탐지 대상 130건을 27설정에서 재집계하면 항상 검출 72, 설정 의존 57, 항상 미검출 1이다. 따라서 기본 설정의 미검출 등을 합친 58건 전체를 ``설정에 따라 바뀜''이라고 서술하지 않는다. 반복 판단의 시계를 처음부터 보존할지 재시작할지에 따라 두 장면의 $D=3$초 판정이 달랐고, 연결 근거 없는 장면 경계의 요구 이월도 별도 시험했다. 이런 결과는 설정과 자료 연결 근거를 보존해야 함을 보여 주며 차량 안전 기준을 확립하지 않는다.
\end{revision}

\section{논의}
\label{sec:discussion}

\begin{revision}

\subsection{검사 가능성과 도구의 역할}
조건별 의무와 해제 증거를 연결하면, 새로운 근거 B가 앞선 A를 해제하는지와 A가 미해제인지 미상인지를 구분할 수 있다. 최초 문제 사건ㆍ관련 의무ㆍ사유를 원문 구절에 연결하는 출력은 수정할 주석이나 추가 증거가 필요한 항목을 검토하는 데 쓸 수 있다. 그러나 작성 프레임의 정확성을 모델 검사가 보장하지는 않는다. 잘못 연결한 조건 ID나 누락된 의무는 별도 의미 검토가 필요하다.

이전 이력 창과 전체 이력의 비교는 기억 길이가 일부 판정에 영향을 주었음을 보여 준다. FSM의 144/144와 보완 모델의 동일 사례 144/144는 이 범위의 논리를 단순 상태 기계로도 실행할 수 있다는 근거다. UPPAAL의 역할은 표현한 상태ㆍ조건 속성을 명시적으로 질의하고 진단 경로와 모델 파일을 보존하는 데 있다. 정확도 우위나 UPPAAL만의 필요성을 입증하지 않는다.

\subsection{주석 검토에서 구분할 조치}
확정 모순은 미해제가 확인된 의무와 행동을 함께 검토한다. 근거 부족은 A의 해제 여부를 B로 대신 채우지 않고 추가 주석/관측을 요청할 대상으로 둔다. 양립 판정도 명시된 의무에 대해서만 유효하다. P4는 주석을 수정할 근거가 아니라 검사 상태 처리의 고장을 점검할 근거다. 따라서 모든 경고를 자연 오류라고 하거나 자료를 자동 삭제하는 절차로 해석하지 않는다.

자연 자료의 낮은 일치도와 합의 사유 누락은 현재 기준을 실제 오류 정답으로 쓰기 어렵게 한다. 이번 작성 시나리오는 이 문제를 해결한 사람이 평가한 gold가 아니다. 자동 추출의 전 사례 미상 결과도 별도 한계로 남는다. 독립 자료와 실제 전문가 검토는 후속 검증이며, 현재의 적용 가능성 실증을 그런 성능으로 확대하지 않는다.
\end{revision}

\section{한계와 결과 해석 시 주의점}
\label{sec:threats}

\begin{revision}

\subsection{정답과 독립성}
기존 40쌍은 기대 결과를 확인해 남긴 규칙 회귀시험이다. 새 144건은 최초 결과에 따라 선별하지 않았지만 원문ㆍ프레임ㆍ기대 라벨과 비교기를 같은 설계자가 만들었다. 48개 실패를 사용한 조건별 모델은 같은 사례로 재시험했다. 개발 미사용 독립 자료, 외부 사람 의미 정답과 자연 오류 평가를 확보하지 못했다. 따라서 완전 검출, 독립 검증 완료나 자연 precision/recall/F1을 주장하지 않는다.

\subsection{자동 추출과 의미 표현}
조건별 모델은 자동 NLP가 아니다. 원문의 증거를 조건 ID로 올바르게 연결하는 단계가 작성 입력에 의존한다. 기존 파서는 제한된 영어 문구ㆍ명시적 연결사를 처리하고 관측 증거를 미상으로 둔다. 부정, 암시, 지시어, 복합행동, 복수 의무와 다른 언어에 대한 일반적 의미 추출은 검증하지 않았다. 단일 의무 투영 손실과 대기 중 누적 미상이 최초 결과를 낮췄다. 보완 모델은 모든 P1--P4 고장이나 임의의 다중 선행조건을 지원한 모델이 아니다.

\subsection{증거의 유효성과 모델 범위}
증거 보존은 같은 대상ㆍ조건과 명시한 구간이 유지된다는 가정을 쓴다. 객체 추적, 관측 노후화ㆍ만료, 해제 후 재등장과 시간에 따라 변하는 유효성은 포함하지 않았다. 복수 동시 문제의 사유 플래그를 보존하지만 대표 최초 의무는 정의한 우선순위로 하나를 선택한다. 이 출력이 모든 의무의 완전한 인과 설명은 아니다. 중심 모델은 이산 사건 순서이며 정량 시간 제한ㆍ브레이크/조향 명령ㆍ차량 동역학ㆍ폐루프 상호작용이 없다. 사건 끝에서 미래 해제를 못 봤다는 사실만으로 위반을 선언하지 않는다.

\subsection{자연 자료와 재현성}
27장면은 행동 표현으로 고른 비무작위 표본이다. 텍스트 불일치 23/27과 합의 사유 메모 0/27이 남아 있고, 검토자 전문성ㆍ교육ㆍ중재 및 논의 절차를 충분히 복원하지 못했다. 사람 재평가를 실시하지 않았다. 원 생성 실행별 모델 revision, prompt/config와 품질 점수 상류 임계값이 미복원이라 완전 재생성을 주장할 수 없다. 현재 입력ㆍ코드ㆍXMLㆍ질의ㆍ로그ㆍ해시는 보존했지만 restricted 자료의 외부 공개나 독립 사용 이력을 인증한 것은 아니다.

\subsection{보조 반응 분석과 일반화}
속도 변화는 실제 제어 명령의 대용이며 객체 해제의 참ㆍ거짓을 제공하지 않는다. $W,\delta,\rho,D$ 및 이미 정지 규칙은 탐색적 설정이고 안전 법규에서 유도하지 않았다. 27설정은 가능한 모든 조건을 포함하지 않는다. 궤적 불일치 7건을 자연 CoC 모순 정답으로 바꿀 수 없다. 다른 자료ㆍ언어ㆍ생성 모델로의 일반화, 실제 차량 안전과 학습 성능 개선은 미검증이다.
\end{revision}

\section{결론}
\label{sec:conclusion}

\begin{revision}

본 논문은 연속 CoC의 의무와 해제 조건을 후속 증거에 연결하고, 사건 순서에서 모순과 근거 부족을 구분하여 최초 문제 사건ㆍ의무ㆍ사유를 추적하는 프레임워크의 적용 가능성을 다루었다. 새 행동의 근거 B와 앞선 의무 A의 해제 증거를 구분하고, A 미해제 확인과 해제 여부 미상을 서로 다른 결과로 기록한다.

144개 작성 시나리오의 최초 시험에서 기존 구조화 경로는 96/144, 최초 위반은 54/72였다. 정상 36건 중 30건은 미상으로 남았고 자동 텍스트 경로는 전 사례 미상이었다. 단순 FSM은 144/144였다. 최초 실패 48건을 분석한 후 별도 조건별 UPPAAL 모델은 같은 사례에서 판정 144/144, 최초 위반 72/72, 문제 위치ㆍ의무ㆍ사유 각 108/108을 보였다. 이는 개발 후 동일 사례 재시험이며 독립 일반화 성능이 아니다.

기존 Python--UPPAAL의 판정ㆍ오류 종류ㆍ최초 위치는 144건에서 일치했다. 반례 접두 경로의 미상 사유 차이 24건은 종료 질의로 진단했으며 원 파서는 변경하지 않았다. 자연 검토는 309인접쌍 목록 중 별도 선정한 27장면으로, 텍스트 $\kappa=0.027301$과 합의 사유 기록 부재 때문에 자연 오류 정답으로 쓰지 않았다. 기존 P4 10쌍은 주석 오류와 분리한 처리 회귀시험이고 보조 반응 분석도 차량 안전 성능으로 해석하지 않는다.

작성된 미해제 충돌 사례에서는 이전 의무의 보존이 해당 사건 관계를 판정하는 데 필요했다. 다만 이 관찰을 모든 CoC에 대한 필요성ㆍ효과나 UPPAAL의 정확도 우위로 일반화하지 않는다. 현재 결과는 명시된 프레임의 시간ㆍ조건 관계를 실행하고 검토 근거를 추적하는 방법 실증이다. 자동 의미 추출, 개발 미사용 독립 시험, 사람 판정 신뢰도, 자연 오류 성능과 물리적 안전은 후속 검증으로 남는다.
\end{revision}


\section*{감사의 글}
\begingroup
\raggedright
본 연구는 대한민국 정부(교육과학기술부)의 재원으로 한국연구재단(NRF)의 지원을 받아 수행되었으며(\nolinkurl{NRF-2023R1A2C1006639}), 대한민국 정부(과학기술정보통신부)의 재원으로 정보통신기획평가원(IITP)의 지역지능화혁신인재양성사업 지원을 받아 수행되었음(\nolinkurl{IITP-2026-RS-2024-00436773}).
\par
\endgroup
\bibliographystyle{unsrt}
\begingroup
\fontsize{8pt}{12pt}\selectfont
\bibliography{references}
\endgroup
\end{document}
